\originalchapter{对称性与离散波}\label{ch:chain}
\originalsection{对称变换为何帮助求模}
两个相同摆交换位置以后, 运动定律不变. 数学上, 若线性变换 $S$ 满足 $S(M^{-1}K)=(M^{-1}K)S$, 则对任一解 $\boldsymbol x(t)$, $S\boldsymbol x(t)$ 仍是解. $[A,B]=AB-BA$ 称为交换子; $[S,M^{-1}K]=0$ 表示对称操作与动力学相容.

若一个频率对应一维本征子空间, $S$ 将该子空间映到自身, 因而此模也是 $S$ 的本征向量. 反过来, 若 $S$ 的某个本征子空间是一维, 与它交换的动力学矩阵只能在其中乘一个数. 对交换两个摆的 $P$, $P^2=I$, 特征值为 $\pm1$, 向量分别为 $(1,1)$ 与 $(1,-1)$. 若使用空间镜像, 位移方向也反号, 此时 $S=-P$, 特征值标签相反, 物理模形状不变. 这个结论依赖两个部件相同及边界对称, 并不适用于任意双自由度系统. 简并子空间内仍需进一步解动力学.

\originalsection{无限链的传播与色散}
一列相同质量 $m$ 的物块, 平衡间距 $a$, 最近邻用劲度 $K_s$ 的弹簧连接. $u_j$ 是第 $j$ 个物块沿链的位移, 不包含平衡坐标 $ja$. 左右弹簧的伸长差给
\begin{equation}\label{eq:discrete}
m\ddot u_j=K_s(u_{j-1}-2u_j+u_{j+1}).
\end{equation}
平移一个格点不改变方程. 设平移本征关系 $a_{j+1}=\beta a_j$, 则 $a_j=a_0\beta^j$. 对向左右无限延伸且幅度有界的振动, 取 $|\beta|=1$, 写 $\beta=\e^{\ii ka}$; $k$ 是波数, 单位 $\mathrm{m}^{-1}$. 代入 $u_j=\Rea[A\e^{\ii(kja-\omega t)}]$ 得
\begin{equation}\label{eq:chain-dispersion}
\omega^2(k)=\frac{2K_s}{m}(1-\cos ka)
=\frac{4K_s}{m}\sin^2\frac{ka}{2}.
\end{equation}
这就是色散关系, 说明不同空间变化尺度对应什么振动频率. $k$ 与 $k+2\pi/a$ 给相同格点值, 所以只需取第一布里渊区 $-\pi/a\le k\le\pi/a$. 正频支可写 $\omega=2\sqrt{K_s/m}|\sin(ka/2)|$. $k=0$ 为整体平移; $|k|=\pi/a$ 时相邻物块反相, 频率达到最大 $2\sqrt{K_s/m}$. 选取 $\pm k$ 的线性组合可形成实的正弦、余弦驻波, 选取随时间的适当组合可形成行波. 空间本征模不一定只能解释为驻波.

\originalsection{边界选择允许的波数}
将无限链的解限制到有限链, 必须同时满足边界. 先写 $a_j=C\e^{\ii jka}+D\e^{-\ii jka}$. 对 $N$ 个物块和固定端 $u_0=u_{N+1}=0$, 空间振幅为
\[
a_j^{(n)}=\sin\frac{n\pi j}{N+1},\qquad
k_na=\frac{n\pi}{N+1},\qquad n=1,\ldots,N.
\]
第一个边界令正、负波数的系数互为相反数, 第二个要求 $\sin[(N+1)ka]=0$. $n=0$ 给零向量, $n=N+1$ 也给零向量, 都不是正常模. 对 $N$ 个自由端物块, 用虚拟格点条件 $u_0=u_1$, $u_{N+1}=u_N$ 描述端点只有一根弹簧. 第一条件给 $D=C\e^{\ii ka}$, 空间振幅遂正比于 $\cos[(j-1/2)ka]$. 第二条件给 $\sin(Nka)\sin(ka/2)=0$, 因而
\[
a_j^{(n)}=\cos\left[\frac{n\pi}{N}\left(j-\frac12\right)\right],\quad
k_na=\frac{n\pi}{N},\quad n=0,\ldots,N-1.
\]
这里 $n=0$ 是允许的零频平移模; $n=N$ 给全零向量, 必须排除. 两组边界不能共用同一个模数范围.

固定链的离散正交关系为 $\sum_{j=1}^N\sin(n\pi j/(N+1))\sin(r\pi j/(N+1))=(N+1)\delta_{nr}/2$. 有限三角级数可直接验证. 给定初位移 $u_j^0$ 与初速 $v_j^0$, 模态解 $u_j=\sum_n[A_n\cos\omega_nt+B_n\sin\omega_nt]a_j^{(n)}$ 的系数为
\[
A_n=\frac2{N+1}\sum_{j=1}^Nu_j^0a_j^{(n)},\qquad
B_n=\frac2{(N+1)\omega_n}\sum_{j=1}^Nv_j^0a_j^{(n)}.
\]
自由链余弦模的非零模范数平方为 $N/2$, 按同样方式投影; 零模另外给 $\bar u^0+\bar v^0t$, 其中横线是全部物块初值的算术平均. 从边界得到模式, 从初值得到每个模式占多少, 是两个不同问题.

若左端固定, 右端以 $u_{N+1}=\Delta\cos\omega_dt$ 驱动, 在分母非零时特解可写
\[
u_j(t)=\Delta\frac{\sin(jk_da)}{\sin[(N+1)k_da]}\cos\omega_dt,
\qquad \omega_d^2=\frac{4K_s}{m}\sin^2(k_da/2).
\]
$k_da=0$ 和 $\pi$ 使上式成为可去的 $0/0$, 分别给静态极限 $u_j=\Delta j/(N+1)$ 和带边响应 $u_j=\Delta(-1)^{j-N-1}j\cos\omega_dt/(N+1)$, 可用正弦比值极限或直接代入差分式验证. 它们不是有限链的正常模共振. 真正共振的零点为 $k_da=n\pi/(N+1)$, $1\le n\le N$. 当外界耦合到对应模时, 必须包含增长或阻尼, 不能继续用此稳态公式. 驱动频率超过链的最高频率时, 没有实 $k_d$; 复波数给相邻振幅交替并指数衰减的响应, 不能用实反余弦强行计算传播波.

\originalsection{张紧弦的连续极限}
另一种离散链是用张力 $T$ 的轻绳连接点质量, 质量只能横向运动. 设横位移为 $y_j$, 要求相邻绳段斜率 $|(y_{j+1}-y_j)/a|\ll1$; 仅要求绝对位移小并不充分. 小斜率下竖直力为 $T(y_{j-1}-2y_j+y_{j+1})/a$. 若 $\mu=m/a$ 是线密度, 令间距趋小而保持 $\mu,T$ 有限, Taylor 展开给
\[
y(x+a)-2y(x)+y(x-a)=a^2y_{xx}(x)+O(a^4),\qquad
 y_{tt}=c^2y_{xx},\quad c=\sqrt{T/\mu}.
\]
连续近似要求 $|ka|\ll1$, 即波长远大于格点间距. 离散关系 $\omega^2=4T\sin^2(ka/2)/(ma)$ 在此极限变成 $\omega^2=c^2k^2$. 线密度的单位是 $\mathrm{kg/m}$, 所以 $T/\mu$ 的单位为速度平方. 这个极限建立了弦波方程的物理来源; 不是因为看见二阶偏微分方程就把数学中的任意函数解释为一根真实弦.
